\chapter{Abstract}

Economics may be seen as the science of aligning a complex, intelligent system (the ``market'') to some social goal (e.g. consumer preferences). AI alignment is the problem of ensuring that AI agents pursue their principals' goals. This thesis argues this is not merely a suggestive analogy, but the basis for a productive research program. We develop the equivalence at three levels. 
\begin{enumerate}
\item \textbf{Agent foundations}. We study markets as a model for boundedly-rational agency. Here we first investigate an established line of research where a ``program market'' learns to solve MDPs or price logical sentences; and next propose a graphical economic model that we use to prove long-sought equivalences between (chain) economies with backpropagation, and between (pass-through/production-less) economies and MDPs. Backpropagation on this economy realizes a generalized Bellman recursion, which we call a ``rational theory of reward'' (complementing Bayes, which gives a rational theory for how to update in response to reward/evidence).
\item \textbf{Market failures}. We cast ``Outer Alignment'' as the economic problem of ``Information Asymmetry'', and use this to motivate a scalable oversight protocol we call ``Recursive Information Markets''. In our study of scalable oversight, we found flaws in the existing empirical research on Debate, and introduce a principled approach to empirically evaluating scalable oversight mechanisms based on the \emph{Agent Score Difference} metric: we use this to establish the \emph{Scalable Oversight Mechanism benchmark}.
\item \textbf{Eliciting latent knowledge}. Aligning and monitoring superhuman agents may often involve eliciting agent beliefs on tasks and sentences where we lack ``ground truth'': either because this is specifically the domain where scoring is non-trivial, or because it is relevant for interpretability. We study two approaches: first, we introduce a scoring mechanism for sentences that are \emph{neither verifiable nor falsifiable} i.e. first-order logic sentences that are higher than Level-1 on the arithmetical hierarchy. Second, we develop a consistency-based evaluation pipeline and principled, arbitrage-based inconsistency metrics for language model forecasters (which we then study empirically).
\end{enumerate}
